
Même
dans le cas où ces dernières se ramènent aux classes de Chern des Modules localement libres ($X$ schéma avec Module inversible ample, ou espace compact), la construction de ces classes de Chern ``à supports'' n'est pas faite, et ne pourra sans doute se faire qu'en même temps que celle des classes de Chern générales postulées dans 7.1\SGAAuthorFootnote{*}{Ajouté en octobre 1968. Les problèmes soulevés ici ont été résolus dans le cadre de la \uline{cohomologie de Hodge} par L. ILLUSIE; cf. [13]}

Il est probable que les classes de Chern avec supports auront un rôle à jouer dans la démonstration de la formule de Riemann--Roch cohomologique sans dénominateurs (cf. 3.5). D'autre part, elles devraient permettre de donner une version améliorée de la formule de Riemann--Roch (avec ou sans dénominateurs) en introduisant dans cette formule des familles de supports.

\section*{8. Anneau de Chow des schémas réguliers}

Le problème est de définir l'équivalence rationnelle des cycles, et d'étudier l'anneau de classes obtenu (notamment ses propriétés fonctorielles), pour un schéma noethérien régulier $X$ admettant un Module inversible ample, en généralisant la théorie connue lorsque $X$ est lisse sur un corps [4]. Le résultat technique clef sera encore un ``moving lemma'', assurant qu'on peut faire varier un cycle dans sa classe de façon à intersecter avec la bonne dimension un cycle donné.

Le problème envisagé est assez naturel à priori, d'autant plus que nous avons vu (4.5 à 4.8) que les substituts $\operatorname{Gr}^\bullet(X)$, $\Gtop^\bullet(X)$ de l'anneau de Chow $A(X)$ ne sont guère satisfaisants quand on ne se borne pas à une théorie ``modulo torsion''. Lorsque $X$ se déduit d'un schéma algébrique projectif et lisse $Y$ sur un corps $k$, en prenant le schéma local au sommet d'un cône projetant $C$ et enlevant l'origine, on trouve (en admettant que $A(X)$ est bien la limite inductive des anneaux de Chow des $U-(O)$, où $U$ parcourt les voisinages ouverts dans $C$ du sommet $O$) que $A(X)$ est canoniquement isomorphe à l'anneau quotient $A(Y)/\xi A(Y)$, où $\xi\in A^1(Y)$ est induit par l'immersion projective considérée de $Y$. Cela montre que la connaissance de l'anneau de Chow $A(X)$ (de type non classique) équivaut à la connaissance de $A(Y)/\xi A(Y)$, qu'on pourrait appeler la ``partie primitive'' de l'anneau de Chow $A(Y)$ (en un sens suggéré par la théorie cohomologique de Lefschetz--Hodge des variétés algébriques projectives non singulières). Cet exemple montre que les ``anneaux de Chow locaux'' des anneaux locaux à singularité
